\documentclass[11pt]{article}
\usepackage[margin=1in]{geometry}
\usepackage{amsmath,amssymb,amsthm}
\usepackage{hyperref}
\usepackage{booktabs}
\usepackage{enumitem}

\newtheorem{theorem}{Theorem}
\newtheorem{lemma}[theorem]{Lemma}
\newtheorem{definition}[theorem]{Definition}

\title{\textbf{Ground Truth: RSA Correctness from First Principles\\in Pure Lean 4 Without Mathlib}}
\author{Davis Brief\\\texttt{davis@team8.co}}
\date{April 2026}

\begin{document}
\maketitle

\begin{abstract}
We present a complete, machine-checked proof of RSA encryption correctness in the Lean 4 proof assistant, built entirely from first principles without any dependency on the Mathlib library. The development comprises 196 lemmas organized across 12 tiers, beginning with basic natural number arithmetic ($0 + n = n$) and culminating in the theorem that RSA decryption inverts encryption ($m^{de} \equiv m \pmod{pq}$ for appropriate primes $p, q$ and exponents $d, e$). The proof chain passes through divisibility theory, the Euclidean algorithm, prime factorization, Fermat's Little Theorem (proved via binomial coefficients), B\'{e}zout's identity, and the Chinese Remainder Theorem. The entire development compiles with zero \texttt{sorry} axioms and depends only on Lean's three core axioms (\texttt{propext}, \texttt{Quot.sound}, \texttt{Classical.choice}). To our knowledge, this is the first standalone formalization of RSA correctness in Lean 4 without Mathlib. The proof artifact is publicly verifiable: any user can clone the repository and confirm correctness by running \texttt{lake build}.
\end{abstract}

\section{Introduction}

RSA~\cite{rsa1978} is the most widely deployed public-key cryptosystem in the world. Its correctness relies on a chain of number-theoretic results stretching from basic arithmetic through Fermat's Little Theorem and the Chinese Remainder Theorem. While these results are well-known and have been formalized in various proof assistants (notably in Mathlib for Lean 4~\cite{mathlib}, and in Coq, Isabelle/HOL, and others), every existing Lean 4 formalization depends on the Mathlib library, a 1.5-million-line dependency that, while powerful, obscures the logical structure of the proof.

We present \textbf{Ground Truth}, a self-contained formalization of RSA correctness in 2,553 lines of pure Lean 4 code with zero external dependencies beyond the language's standard library. The development is organized as a linear chain of 12 ``tiers,'' each importing only its predecessor, so that the logical dependency structure is immediately apparent.

Our contributions are:
\begin{enumerate}[leftmargin=*]
    \item A complete, dependency-free proof of RSA correctness in Lean 4, comprising 196 machine-checked lemmas with zero \texttt{sorry}.
    \item A proof of Fermat's Little Theorem from scratch, via binomial coefficients and the ``freshman's dream'' identity, without any finite field or group theory machinery.
    \item A constructive proof of B\'{e}zout's identity and the existence of modular inverses, using an integer-valued extended GCD.
    \item A demonstration that a small LLM (Gemma 4 E4B, 4.5B parameters; DeepSeek-Prover-V2 7B) can contribute meaningfully to formal proof development when paired with Lean's type-checker as an oracle.
\end{enumerate}

\section{Architecture}

The proof is organized into 12 tiers, each building strictly on the previous:

\begin{table}[h]
\centering
\begin{tabular}{@{}clcc@{}}
\toprule
\textbf{Tier} & \textbf{Topic} & \textbf{Lemmas} & \textbf{Key Result} \\
\midrule
1  & Natural number basics     & 20 & $a + b = b + a$ \\
2  & Ordering                  & 20 & Trichotomy \\
3  & Division \& modular arithmetic & 19 & $a = b \cdot \lfloor a/b \rfloor + a \bmod b$ \\
4  & Divisibility              & 19 & $d \mid a \wedge d \mid b \Rightarrow d \mid (a+b)$ \\
5  & Lists \& structural induction & 19 & $|l_1 \mathbin{++} l_2| = |l_1| + |l_2|$ \\
6  & GCD \& coprimality        & 18 & Euclid's lemma for coprime pairs \\
7  & Primes                    & 14 & $\exists$ infinitely many primes \\
8  & Fermat's Little Theorem   & 27 & $a^p \equiv a \pmod{p}$ \\
9  & B\'{e}zout \& modular inverse & 12 & $\gcd(a,n)=1 \Rightarrow \exists b.\; ab \equiv 1 \pmod{n}$ \\
10 & Chinese Remainder Theorem & 12 & CRT existence \& uniqueness \\
11 & Fermat-to-RSA bridge      & 10 & $m^{de} \equiv m \pmod{p}$ \\
12 & RSA correctness           &  6 & $m^{de} \equiv m \pmod{pq}$ \\
\midrule
   & \textbf{Total}            & \textbf{196} & \\
\bottomrule
\end{tabular}
\caption{Proof chain from basic arithmetic to RSA correctness.}
\label{tab:tiers}
\end{table}

Each tier is a single Lean 4 file. Every lemma carries a plain-English comment explaining what it proves, making the development accessible to non-experts.

\section{Key Proofs}

\subsection{Fermat's Little Theorem (Tier 8)}

The proof of $a^p \equiv a \pmod{p}$ for prime $p$ proceeds in three stages:

\paragraph{Stage 1: Binomial coefficient infrastructure.} We define binomial coefficients via Pascal's triangle:
\[
C(n, 0) = 1, \quad C(0, k+1) = 0, \quad C(n+1, k+1) = C(n,k) + C(n, k+1)
\]
and prove the product identity $C(n,k) \cdot k! \cdot (n-k)! = n!$ by induction on $n$, with a case split on whether $k < n$ or $k = n$.

\paragraph{Stage 2: Prime divides middle binomial coefficients.} Using the product identity and the fact that $p$ does not divide $k!$ or $(p-k)!$ for $0 < k < p$ (proved by showing a prime cannot divide a factorial of smaller numbers), we establish:
\[
0 < k < p \implies p \mid C(p, k)
\]

\paragraph{Stage 3: Freshman's dream and induction.} We define a partial binomial sum $\text{binomSum}(a, n, j) = \sum_{i=0}^{j-1} C(n,i) \cdot a^i$, prove the binomial theorem $(a+1)^n = \text{binomSum}(a, n, n)$, and show that all middle terms vanish mod $p$. This yields the ``freshman's dream'':
\[
(a+1)^p \equiv a^p + 1 \pmod{p}
\]
Fermat's Little Theorem then follows by induction on $a$.

\subsection{Chinese Remainder Theorem (Tier 10)}

The CRT proof is constructive: given coprime $m, n$ and target remainders $r_1, r_2$, we compute modular inverses of $n \bmod m$ and $m \bmod n$ (via B\'{e}zout's identity from Tier 9), then construct the explicit solution $x = r_1 \cdot n \cdot n^{-1} + r_2 \cdot m \cdot m^{-1} \pmod{mn}$.

Uniqueness follows from the key lemma: if $a \equiv b \pmod{m}$ and $a \equiv b \pmod{n}$ with $\gcd(m,n) = 1$, then $mn \mid (a - b)$, hence $a \equiv b \pmod{mn}$.

\subsection{RSA Correctness (Tier 12)}

The capstone theorem states: for distinct primes $p, q$, if $de \equiv 1 \pmod{(p-1)(q-1)}$, then
\[
m^{de} \equiv m \pmod{pq}
\]

The proof proceeds in three steps:
\begin{enumerate}
    \item \textbf{Mod $p$}: Case split on whether $p \mid m$. If yes, both sides are $0 \bmod p$. If no, apply the multiplicative form of Fermat's Little Theorem: $m^{de} = m \cdot m^{k(p-1)(q-1)} \equiv m \cdot 1 \equiv m \pmod{p}$.
    \item \textbf{Mod $q$}: Symmetric argument.
    \item \textbf{Combine via CRT}: Since $\gcd(p,q) = 1$ (distinct primes are coprime), the Chinese Remainder Theorem yields $m^{de} \equiv m \pmod{pq}$.
\end{enumerate}

\section{Proof Methodology}

The development was produced through a human-AI collaboration over a single session. The workflow:

\begin{enumerate}[leftmargin=*]
    \item \textbf{Scaffolding}: All 196 theorem statements were written by hand with \texttt{sorry} placeholders and verified to compile.
    \item \textbf{Automated tactics}: A Python orchestrator (\texttt{proof\_farm.py}) tested each lemma against a battery of 40+ tactic combinations (\texttt{omega}, \texttt{simp}, \texttt{induction}, \texttt{cases}, etc.). This solved approximately 60\% of lemmas instantly.
    \item \textbf{LLM-assisted proving}: Lemmas that resisted automated tactics were sent to Gemma 4 E4B (4.5B parameters) and DeepSeek-Prover-V2 (7B parameters) running on a local RTX 2060. Lean's type-checker served as the oracle; candidates were accepted only if they compiled.
    \item \textbf{Human intervention}: Hard lemmas (Fermat's Little Theorem, B\'{e}zout's identity, CRT, prime existence) were proved by the human author with AI assistance for tactic suggestions and arithmetic rearrangements.
\end{enumerate}

DeepSeek-Prover-V2 proved several lemmas on first attempt (e.g., \texttt{mod\_lt} via \texttt{exact Nat.mod\_lt \_ hb}), while Gemma 4 E4B contributed primarily through trial-and-error on simple arithmetic goals. Neither model could handle proofs requiring novel multi-step reasoning about custom definitions.

\section{Axioms and Trust Base}

The development depends on exactly three axioms, all part of Lean 4's core:
\begin{itemize}
    \item \texttt{propext}: Propositional extensionality
    \item \texttt{Quot.sound}: Quotient soundness
    \item \texttt{Classical.choice}: The axiom of choice (used via \texttt{by\_cases} for decidability)
\end{itemize}

No custom axioms are introduced. The \texttt{sorry} count across all 12 tier files is exactly zero.

\section{Related Work}

RSA correctness has been formalized in several proof assistants:
\begin{itemize}
    \item \textbf{Coq}: Bertot and Cast\'{e}ran~\cite{bertot2004} include RSA as an exercise in \textit{Interactive Theorem Proving and Program Development}.
    \item \textbf{Isabelle/HOL}: The AFP contains formalizations of RSA-related number theory.
    \item \textbf{Lean 4 + Mathlib}: Mathlib's \texttt{ZMod} module provides the infrastructure for RSA, but no standalone proof exists.
\end{itemize}

To our knowledge, Ground Truth is the first formalization of RSA correctness in Lean 4 that does not depend on Mathlib or any external library.

\section{Conclusion}

We have demonstrated that a complete, machine-checked proof of RSA correctness can be built from first principles in 2,553 lines of pure Lean 4, organized as a readable 12-tier chain from $0 + n = n$ to $\text{decrypt}(\text{encrypt}(m)) = m$. The development is self-contained, dependency-free, and publicly verifiable.

The proof artifact is available at: \texttt{[repository URL upon publication]}

\begin{thebibliography}{9}
\bibitem{rsa1978}
R.~Rivest, A.~Shamir, and L.~Adleman.
\newblock A method for obtaining digital signatures and public-key cryptosystems.
\newblock \textit{Communications of the ACM}, 21(2):120--126, 1978.

\bibitem{mathlib}
The mathlib Community.
\newblock The {L}ean mathematical library.
\newblock In \textit{CPP 2020}, pages 367--381, 2020.

\bibitem{bertot2004}
Y.~Bertot and P.~Cast\'{e}ran.
\newblock \textit{Interactive Theorem Proving and Program Development: Coq'Art}.
\newblock Springer, 2004.
\end{thebibliography}

\end{document}
